% !PRINTED 404
\weberpage{404}

The question to which we are led here is thus how one can
reduce the solution of an algebraic equation, whose coefficients
depend on a certain number of variables, to an equation
with the smallest possible number of parameters, and in
particular, under what circumstances and by what means
one can reduce it to equations with only one
parameter.

We accordingly now pose the question, under what
assumptions, for a \gesperrt{general equation of degree $n$,
resolvents with only one parameter exist.}

Let $x_0, x_1, \ldots, x_{n-1}$ be a system of $n$ independent variables,
\weq{1}{u = \Phi(x_0, x_1, \ldots, x_{n-1})}
a rational function of these variables which, under the permutations
of the \gesperrt{alternating group} of the $n$ letters $x$, passes over into the
mutually distinct functions
\weq{2}{u,\ u_1, u_2, \ldots u_{\nu-1}}
where $\nu$ may be equal to or less than the degree of the alternating
group, and $z$ is a function of the variables $x$
which remains unchanged under the alternating group. The
question is when a rational equation of degree $\nu$ with respect
to $u$ exists,
\weq{3}{F(u, z) = 0,}
which is identically satisfied by the $\nu$ functions (2)? The
coefficients in (3) must be independent of the $x$ and hence pure
numbers. For the present we let no restriction of the domain of rationality
enter in the region of numbers.

Equation (3) is, when $\nu > 1$, a resolvent of the
equation of degree $n$ whose roots are the $x$, which depends only on the
\gesperrt{one parameter} $z$. We exclude the case $\nu = 1$.

\begin{center}
\S.\,107.

\gesperrt{Theorem of L\"uroth.}
\end{center}

First the following lemma is to be proved \textsuperscript{1}):

1. \gesperrt{If $u, u_1, u_2, \ldots, u_{\nu-1}$ is a system of rational functions
 of one variable $t$, then there is a rational
function of the $u, u_1, \ldots, u_{\nu-1}$:}
\[
\vartheta = \chi(u, u_1, \ldots, u_{\nu-1})
\]

\begin{center}
\rule{0.3\linewidth}{0.4pt}
\end{center}

\textsuperscript{1}) \gesperrt{L\"uroth}, Mathematische Annalen, Bd. IX.
%CONT>>
